%When it comes to making relevant recommendations, time has always been an important factor for many reasons. For example, the characteristics of an item (e.g. popularity or relevance) can vary dramatically along with the time shift. Users are also heavily influenced by time since their interests continuously change and evolve over time. Furthermore, the temporal distance between the two interactions can tell a lot about their relationship. For these reasons, many works have tried to utilize time information effectively.

%TimeSVD++\cite{TIMESVD++}, which included time factor into the traditional matrix factorization method, was one of the main contributions for the winning of Netfliz Grand Prize \cite{NETFLIXPRIZE}. Bayesian Probabilistic Tensor Factorization (BPTF)\cite{BPTF} added time constraints when factorizing the tensors. However, these works were limited to the traditional factorization settings where the sequential information is ignored. Meanwhile, Time-LSTM\cite{TIMELSTM} equipped LSTMs with several forms of time gates to better model the time intervals in user's interaction sequence. Recently, CTA\cite{CTA} used temporal information crucially by using multiple parametrized kernel functions to calibrate the self-attention mechanism.


Since temporal information hold contextual information, they can be very crucial for recommendation performance. TimeSVD++ and BPTF~\cite{TIMESVD++, BPTF} adopted time factor into the matrix factorization method, and TimeSVD++ was one of the main contributions for the winning of Netflix Grand Prize~\cite{NETFLIXPRIZE}.
%Bayesian Probabilistic Tensor Factorization (BPTF)~\cite{BPTF} added time constraints when factorizing the tensors. However, these works were limited to the traditional factorization setting where the sequential nature of data is ignored. Meanwhile,
Time-LSTM~\cite{TIMELSTM} equipped LSTMs with several forms of time gates to better model the time intervals in user's interaction sequence. Recently, CTA~\cite{CTA} used multiple parametrized kernel functions on temporal information to calibrate the self-attention mechanism.